The goals of this section are to revisit the Betti map, the Betti form and make a link between them. 
In this paper we construct the Betti map using the universal family of
principally polarized abelian varieties with level-$\ell$-structure and
bypass the ad-hoc construction found in \cite{GaoHab}.

 In this section we will make
the following assumptions.
All varieties  are defined over the field $\IC$.
Let $S$ be an  irreducible, regular, quasi-projective variety over $\IC$. 
Let $\pi\colon \cA\rightarrow S$ be an abelian scheme of  relative
dimension $g$, that carries a principal polarization, and such that
 $\cA$ is equipped with level-$\ell$-structure, for some
 $\ell\ge\bestlevel$, \textit{i.e.}, {\tt (Hyp)} is satisfied. 

\begin{proposition}\label{PropBettiMap}
Let $s_0 \in S(\IC)$. Then there exist an open neighborhood $\Delta$ of $s_0$ in $S^{\mathrm{an}}$, and a map $b_{\Delta} \colon \cA_{\Delta} := \pi^{-1}(\Delta) \rightarrow \mathbb T^{2g}$, called the Betti map, with the following properties.
\begin{enumerate}
\item[(i)] For each $s \in \Delta$ the restriction $b_{\Delta}|_{\cA_s(\IC)} \colon \cA_s(\IC) \rightarrow \mathbb T^{2g}$ is a group isomorphism.
\item[(ii)] For each $\xi \in \mathbb T^{2g}$ the preimage $b_{\Delta}^{-1}(\xi)$ is a complex analytic subset of $\cA_{\Delta}$.
\item[(iii)] The product
$(b_{\Delta},\pi) \colon \cA_{\Delta} \rightarrow \mathbb
T^{2g} \times \Delta$ is a real analytic isomorphism.% 
\end{enumerate}
\end{proposition}

The properties (i) -- (iii) do not uniquely determine $b_\Delta$.
Indeed,  composing $b_\Delta$ with an automorphism of the
topological group $\IT^{2g}$, \textit{i.e.}, an element of
$\mathrm{GL}_{2g}(\IZ)$, leads to a new Betti map satisfying (i) --
(iii). After shrinking   $\Delta$ we may assume that it is connected.
In this case, an application of the
Baire Category Theorem shows that $b_\Delta$ is
uniquely determined by (i) and (iii)
up to composition with a unique element of $\mathrm{GL}_{2g}(\IZ)$.


Andr\'{e}, Corvaja, and Zannier~\cite{ACZBetti} recently began the
study of the maximal rank of the Betti map, especially the submersivity, using a slightly different
definition. A full study of this maximal rank was realized in \cite{GaoBettiRank}. Closely related to the
Betti map is the Betti form, a semi-positive $(1,1)$-form on
$\an{\cA}$, which was first introduced in Mok \cite{Mok11Form}.

\begin{proposition}\label{PropBettiForm}
There exists a closed $(1,1)$-form $\omega$ on $\an{\cA}$, called the Betti form, such that the following properties hold.
\begin{enumerate}
\item[(i)] The $(1,1)$-form $\omega$ is semi-positive, \textit{i.e.},  at each point
the associated Hermitian form is positive semi-definite.
\item[(ii)] For all $N\in\IZ$ we have $[N]^*\omega = N^2 \omega$.
\item[(iii)] If $X$ is an irreducible subvariety of $\cA$ of dimension
$d$ and $\Delta\subset \an{S}$ is open with $\sman{X}\cap\cA_\Delta\not=\emptyset$, then
\[
\omega|_{\sman{X}}^{\wedge d} \not\equiv 0 \quad\text{if and only
if}\quad  \max_{x \in
\sman{X} \cap \cA_{\Delta}} \mathrm{rank}_{\IR}
(\mathrm{d}b_{\Delta}|_{X^{\mathrm{sm,an}}})_x = 2d. %\text{ for some non-empty open }\Delta\subseteq S^{\mathrm{an}}.
\]

\end{enumerate}
\end{proposition}


We will prove both propositions during the course of this section using the universal abelian variety. A dynamical approach can be found in \cite[$\mathsection$2]{CGHX:18}.

\subsection{Betti map for the universal abelian variety}\label{SubsectionBettiMap}
Our proof of Proposition~\ref{PropBettiMap} follows the construction in \cite[$\mathsection$3-$\mathsection$4]{GaoBettiRank}. We divide it into several steps.

We start to prove Proposition~\ref{PropBettiMap} for $S = \mathbb
A_g$, the moduli space of principally polarized abelian variety of
dimension $g$ with level-$\ell$-structure;  it is a fine moduli
space. Let $\pi^{\mathrm{univ}} \colon \mathfrak{A}_g \rightarrow \mathbb A_g$ be the universal abelian variety.

The universal covering $\mathfrak{H}_g \rightarrow \an{\mathbb A}_g$,
where $\mathfrak{H}_g$ is the Siegel upper half space, gives a
polarized family of abelian varieties
$\cA_{\mathfrak{H}_g} \rightarrow \mathfrak{H}_g$ fitting into the diagram
\[
\xymatrix{
\cA_{\mathfrak{H}_g} := \mathfrak{A}_g \times_{\an{\mathbb A}_g}\mathfrak{H}_g \ar[r]^-{u_B} \ar[d] & \an{\mathfrak{A}}_g \ar[d]^{\pi^{\mathrm{univ}}} \\
\mathfrak{H}_g \ar[r] & \an{\mathbb A}_g.
}
\]
For the universal covering $u \colon \IC^g \times \mathfrak{H}_g
\rightarrow \cA_{\mathfrak{H}_g}$ and for each $Z \in
\mathfrak{H}_g$, the kernel of $u|_{\mathbb{C}^g \times \{Z\}}$ is
$\mathbb{Z}^g + Z \mathbb{Z}^g$. Thus the map $\IC^g \times
\mathfrak{H}_g \rightarrow \IR^g \times \IR^g \times
\mathfrak{H}_g \rightarrow \IR^{2g}$, where the first map is the
inverse of $(a,b,Z) \mapsto (a + Z  b,Z)$ and the second map is the
natural projection, descends to a \textit{real} analytic map
\[
b^{\mathrm{univ}} \colon \cA_{\mathfrak{H}_g} \rightarrow \mathbb T^{2g}.
\]
 Now for each $s_0 \in \mathbb A_g(\IC)$, there exists a contractible,
 relatively compact, open neighborhood $\Delta$ of $s_0$ in $\mathbb A_g^{\mathrm{an}}$ such that $\mathfrak{A}_{g,\Delta}:=(\pi^{\mathrm{univ}})^{-1}(\Delta)$ can be identified with $\cA_{\mathfrak{H}_g,\Delta'}$ for some open subset $\Delta'$ of $\mathfrak{H}_g$. The composite $b_{\Delta} \colon \mathfrak{A}_{g,\Delta} \cong \cA_{\mathfrak{H}_g,\Delta'} \rightarrow \mathbb T^{2g}$ is  real analytic and satisfies the three properties listed in Proposition~\ref{PropBettiMap}. Thus $b_{\Delta}$ is the desired Betti map in this case. Note that for a fixed (small enough) $\Delta$, there are infinitely choices of $\Delta'$; but for $\Delta$ small enough, if $\Delta_1'$ and $\Delta_2'$ are two such choices, then $\Delta_2' = \alpha \cdot \Delta_1'$ for some $\alpha \in \mathrm{Sp}_{2g}(\IZ) \subset \mathrm{SL}_{2g}(\IZ)$. Thus we have proved Proposition~\ref{PropBettiMap} for $\mathfrak{A}_g \rightarrow \mathbb A_g$.


\subsection{Betti form for the universal abelian variety}

\newcommand{\X}{X}
\newcommand{\Y}{Y}
\newcommand{\rmd}{\mathrm{d}}
\newcommand{\tran}[1]{{{#1}}^{\!^{\intercal}}}


For the universal covering $\mathbf{u} = u_B\circ u \colon \IC^g
\times \mathfrak{H}_g \rightarrow \mathfrak{A}_g^{\mathrm{an}}$, we
will use $(w,Z)$ to denote the coordinates on $\IC^g \times
\mathfrak{H}_g$.
Below $\mathrm{Im}$ denotes imaginary part. 
%% We will decompose an element $Z\in \mathfrak{H}_g$  as $Z = \X
%% + \sqrt{-1}\Y$
%% with $\X$ the real part and $\Y$ the imaginary part of $Z$. 


\begin{lemma}\label{LemmaBettiOnUniversalCovering}
%Under the change of coordinates \eqref{EqChangeOfCoor}, the $2$-form $2\mathrm{d}a \wedge \mathrm{d} b$ becomes
Define
\[
\hat{\omega}^{\mathrm{univ}}:=\sqrt{-1}\partial \overline{\partial} \left( 2 (\mathrm{Im}w)^{\!^{\intercal}} (\mathrm{Im} Z)^{-1} (\mathrm{Im}w) \right).
\]
Then $\hat{\omega}^{\mathrm{univ}}$ is a closed  semi-positive
$(1,1)$-form on $\IC^g\times\mathfrak{H}_g$ satisfying %we have
\begin{equation}
  \label{eq:bettiformaltformula}
  \hat{\omega}^{\mathrm{univ}} = \sqrt{-1} \tran{\left( \rmd Z \Y^{-1}
      \mathrm{Im}(w) - \rmd w\right)} \wedge \Y^{-1} \left( \rmd \overline Z \Y^{-1}
    \mathrm{Im}(w) - \rmd \overline w\right)
\end{equation}
with $\Y = \mathrm{Im}(Z)$; here and below the symbol $\wedge$
is used as a combination of wedge product and matrix multiplication
when appropriate. %\Pcomment{New comment, referee .3}
Moreover, if $N\in\IZ$ and if we denote by $\widetilde{N} \colon \IC^g \times \mathfrak{H}_g \rightarrow \IC^g \times \mathfrak{H}_g$ the map $(w,Z) \mapsto (Nw,Z)$, then $\widetilde{N}^*\hat{\omega}^{\mathrm{univ}} = N^2\hat{\omega}^{\mathrm{univ}}$.
\end{lemma}
\begin{proof}
The $(1,1)$-form $\hat{\omega}^{\mathrm{univ}}$ is  closed since $\mathrm{d} = \partial + \overline{\partial}$. We will prove the semi-positivity by direct computation.

% Using \Pinlinecomment{TODO: Referee .3 wants explanation for formulas
%   on right of next two lines.}
%   \footnote{It is worth explaining the formulae on the right. We hereby do it for $\partial \Y^{-1} = \frac{\sqrt{-1}}{2} \Y^{-1} \mathrm{d}  Z \Y^{-1}$ and the other one is similar. Taking partial derivatives on both sides of $\Y \Y^{-1} = I$, we get $(\partial \Y)\Y^{-1} + \Y \partial Y^{-1} = 0$. So $\partial \Y^{-1} = - \Y^{-1}(\partial \Y)\Y^{-1}$. But $\partial \Y = \partial \mathrm{Im}Z = - \frac{\sqrt{-1}}{2}\mathrm{d} Z$. Hence we get the desired formula for $\partial \Y^{-1}$.}
   
We have the following formulae for partial derivatives
  \begin{equation*}    
%    \label{eq:dwdZ}
    \begin{aligned}
      \overline\partial \mathrm{Im} w &=& \frac{\sqrt{-1}}{2} \mathrm{d}\overline w, \qquad     
      \overline\partial (\Y^{-1}) &=&& -\frac{\sqrt{-1}}{2} \Y^{-1} \mathrm{d} \overline Z
      \Y^{-1},\\
      \partial \mathrm{Im} w &=& -\frac{\sqrt{-1}}{2} \mathrm{d} w,\qquad
      \partial (\Y^{-1}) &=&& \frac{\sqrt{-1}}{2} \Y^{-1} \mathrm{d}  Z \Y^{-1}.
    \end{aligned}
  \end{equation*}
Let us prove the formulae on the right. We hereby do it for $\partial (\Y^{-1}) = \frac{\sqrt{-1}}{2} \Y^{-1} \mathrm{d}  Z \Y^{-1}$ and the other one is similar. Taking partial derivatives on both sides of $\Y \Y^{-1} = I$, we get $(\partial \Y)\Y^{-1} + \Y \partial (Y^{-1}) = 0$. So $\partial (\Y^{-1}) = - \Y^{-1}(\partial \Y)\Y^{-1}$. But $\partial \Y = \partial \mathrm{Im}Z = - \frac{\sqrt{-1}}{2}\mathrm{d} Z$. Hence we get the desired formula for $\partial (\Y^{-1})$.
  
Using these formulae  
and the Leibniz rule (note that $Z = Z^{\!^{\intercal}}$ and hence $\mathrm{d}Z = \mathrm{d}Z^{\!^{\intercal}}$), we get
\begin{align*}
    \hat\omega^{\mathrm{univ}} = \sqrt{-1} \big( &
 (\mathrm{d}w)^{\!^{\intercal}} \Y^{-1} \wedge  \mathrm{d}\overline{w} + (\mathrm{Im}w)^{\!^{\intercal}} \Y^{-1}  \mathrm{d}Z \wedge \Y^{-1}  \mathrm{d}\overline{Z}  \Y^{-1} (\mathrm{Im}w) \\
 & - (\mathrm{Im}w)^{\!^{\intercal}} \Y^{-1} \mathrm{d}Z
 \Y^{-1} \wedge  \mathrm{d}\overline{w} -
 (\mathrm{d}w)^{\!^{\intercal}}   \wedge \Y^{-1}  \mathrm{d}\overline{Z} \Y^{-1} (\mathrm{Im}w) \big). 
\end{align*}
  Rearranging yields the desired equality (\ref{eq:bettiformaltformula}). 
  The associated  form is 
  \begin{equation*}
    H:\bigl((\xi_w,\xi_Z),(\eta_w,\eta_Z)\bigr) \mapsto 
    \tran{\left( \xi_Z \Y^{-1}
        \mathrm{Im}(w) - \xi_w\right)}  \Y^{-1} \left( \overline{\eta_Z} \Y^{-1}
      \mathrm{Im}(w) -  \overline{\eta_w}\right),
  \end{equation*}
  for $\xi_w,\eta_w \in \IC^g$ and $\xi_Z,\eta_Z\in
  \mathrm{Mat}_g(\IC)$ symmetric, is Hermitian and so
  $\hat \omega^{\mathrm{univ}}$ is real. Moreover, 
    \begin{equation*}
    H\bigl((\xi_w,\xi_Z),(\xi_w,\xi_Z)\bigr) = \tran{v} \Y^{-1} \overline v
    \quad\text{with}\quad v =  \xi_Z \Y^{-1}
        \mathrm{Im}(w) - \xi_w.
  \end{equation*}
  But $\Y$ is positive definite as a real symmetric matrix and thus
  positive definite as a Hermitian matrix. We see that $H$
  is positive semi-definite and this implies that
  $\hat\omega^{\mathrm{univ}}$ is positive semi-definite.

  The ``moreover'' part of the lemma is clear.    
\end{proof}

Next we want to show that $\hat{\omega}^{\mathrm{univ}}$ descends to a $(1,1)$-form on $\an{\mathfrak{A}}_g$. To do this, we first show that $\hat{\omega}^{\mathrm{univ}}$ can be written in a simple form under an appropriate change of coordinates.

Define the complex space $\mathcal X_{2g,\mathrm{a}}$, which is the universal covering of $\an{\mathfrak{A}}_g$, as follows:
\begin{itemize}
\item As a real algebraic space, $\mathcal X_{2g,\mathrm{a}} := \mathbb R^{2g} \times \mathfrak{H}_g$.
\item The complex structure on $\mathcal X_{2g,\mathrm{a}}$ is given by
\begin{equation}\label{EqChangeOfCoor}
\mathbb R^{2g} \times \mathfrak{H}_g = \mathbb R^g \times \mathbb R^g \times \mathfrak{H}_g \cong \mathbb C^g \times \mathfrak{H}_g, \qquad (a,b,Z) \mapsto (a+Zb, Z).
\end{equation}
\end{itemize}


\begin{lemma}\label{LemmaBettiFormComplexReal}
Let $\hat{\omega}^{\mathrm{univ}}$ be as in
Lemma~\ref{LemmaBettiOnUniversalCovering}. Then under the change of
coordinates \eqref{EqChangeOfCoor}, we have 
$\hat{\omega}^{\mathrm{univ}}=2\tran{(\mathrm{d}a)} \wedge \mathrm{d}
b$.
\end{lemma}
\begin{proof} For the moment we write $Z = \X + \sqrt{-1}\Y$ with $\X$ and $\Y$ the real and
imaginary part of $Z\in\mathfrak{H}_g$, respectively. 
Note that $w = a + Zb = (a + \X b) + \sqrt{-1}\Y b$. Hence
$\Y^{-1}(\mathrm{Im}w) = b$ and $\mathrm{d}w = \mathrm{d}a +
Z\mathrm{d}b + \mathrm{d}Z  b$. Using this and noting that $Z$
is symmetric, we have that \eqref{eq:bettiformaltformula} becomes
\begin{equation*}
%\label{EqFirstTerm11Form}
\begin{aligned}
    \hat \omega^{\mathrm{univ}} &=\sqrt{-1}\left( \sqrt{-1}\tran{(\rmd
        b)} \Y + \tran{(\rmd b)} \X + \tran{(\rmd a)}\right) \wedge \Y^{-1}
    \left(\rmd a + \X \rmd b - \sqrt{-1} \Y \rmd b\right) \\
    &=\sqrt{-1}\bigl(\sqrt{-1} \tran{(\rmd b)}\wedge\rmd a + \tran{(\rmd b)}\wedge
    \Y \rmd b + \tran{(\rmd b)} \X \wedge \Y^{-1} \rmd a+ \tran{(\rmd
      b)}\X\wedge 
    \Y^{-1}\X \rmd b + \\ 
    &\phantom{=\sqrt{-1}\bigl(}\tran{(\rmd a)} \wedge \Y^{-1} \rmd a 
    + \tran{(\rmd a)}\wedge \Y^{-1}\X\rmd b -
    \sqrt{-1} \tran{(\rmd a)}\wedge \rmd b\bigr).
\end{aligned}
\end{equation*}

 Many terms will vanish. Indeed, if $M$ is a
matrix, then $\tran{(\rmd b)} \wedge M \rmd a = - \tran{(\rmd
a)}\wedge \tran{M} \rmd b $. As $\tran{(\X\Y^{-1})} = \Y^{-1}\X$ and as
$\tran{(\rmd b)} \X \wedge \Y^{-1} \rmd a=\tran{(\rmd
b)} \wedge \X\Y^{-1} \rmd a$ we find
$\tran{(\rmd b)}\X \wedge \Y^{-1} \rmd a + \tran{(\rmd a)}\wedge
\Y^{-1}\X \rmd b=0$. Observe that $\Y$ is symmetric, and so $\tran{(\rmd
b)}\wedge \Y \rmd b = -\tran{(\rmd b)}\wedge \Y\rmd b$ vanishes. Arguing
along the same line and using that
$\Y^{-1}$ and $\X\Y^{-1}\X$ are  symmetric we find $\tran{(\rmd
a)}\wedge \Y^{-1} \rmd a = 0$ and $\tran{(\rmd b)}\X\wedge \Y^{-1}\X=\tran{(\rmd b)}\wedge \X\Y^{-1}\X \rmd
b=0$. We are left with
$    \hat \omega^{\mathrm{univ}} = 2\tran{(\rmd a)} \wedge \rmd b$. 
\end{proof}


\begin{corollary}\label{CorBettiFormVanishingDirection}
  Let $\hat{C}$ be an irreducible, $1$-dimensional, complex analytic
  subset of an open subset of $\mathcal X_{2g,\mathrm{a}}
  = \IR^{2g} \times \mathfrak{H}_g$ and $\sm{\hat {C}}$ its smooth
  locus. Then $\hat{\omega}^{\mathrm{univ}}$ restricted to $\sm{\hat C}$
  is trivial if and only if
  $\hat{C} \subseteq \{r\} \times \mathfrak{H}_g$ for some
  $r \in \mathbb R^{2g}$.
\end{corollary}
\begin{proof}
  First, assume that the coordinates $(a,b)$ of $\IR^{2g}$ are constant on
  $\hat C$. Then $\hat{\omega}^{\mathrm{univ}}$, which is simply
  $2 \tran{(\rmd a)} \wedge \mathrm{d}b$ by
  Lemma~\ref{LemmaBettiFormComplexReal}, vanishes on $\sm{\hat C}$.

  Conversely, suppose that $\hat{\omega}^{\mathrm{univ}}$ vanishes
  identically on $\sm{\hat C}$. This time we use (\ref{eq:bettiformaltformula}) from 
  Lemma~\ref{LemmaBettiOnUniversalCovering}. As $\Y^{-1}$ is positive
  definite we find $\rmd Z \Y^{-1} \mathrm{Im}(w) = \rmd w$ on $\sm{\hat C}$.
  Using the change of coordinates $w = a+Zb$ we deduce $\mathrm{Im}(w) = \Y b$
  and $\rmd w = \rmd a
  + \rmd Z b +Z\rmd b$. So $\rmd Z b  =\rmd Z \Y^{-1} \mathrm{Im}(w)
  = \rmd w=  \rmd a
  + \rmd Z b + Z\rmd b$ on $\sm{\hat C}$. This equality simplifies to
  $\rmd a + Z \rmd b = 0$ on $\sm{\hat C}$. As $a$ and $b$ are real
  valued and as $Z\in \mathfrak{H}_g$ we conclude $\rmd a = \rmd b =0$
  on $\sm{\hat C}$. So $a$ and $b$ are constant on $\hat C$.
\end{proof}


\begin{lemma}\label{LemmaBettiFormUnivAb}
Let $\hat{\omega}^{\mathrm{univ}}$ be as in
Lemma~\ref{LemmaBettiOnUniversalCovering}. Then
$\hat{\omega}^{\mathrm{univ}}$ descends to a  semi-positive
$(1,1)$-form $\omega^{\mathrm{univ}}$ on $\mathfrak{A}_g$. Moreover,
for $N\in\IZ$  we have $[N]^*\omega^{\mathrm{univ}} = N^2 \omega^{\mathrm{univ}}$.
\end{lemma}
\begin{proof}
Let $\mathrm{Sp}_{2g}$ be the symplectic group defined over $\mathbb Q$, and let $V_{2g}$ be the vector group over $\mathbb Q$ of dimension $2g$. Then the natural action of $\mathrm{Sp}_{2g}$ on $V_{2g}$ defines a group $P_{2g,\mathrm{a}}:= V_{2g} \rtimes \mathrm{Sp}_{2g}$.

We use the classical action of $\mathrm{Sp}_{2g}(\mathbb R)$  on
$\mathfrak{H}_g$, it is transitive. The real coordinate on $\mathcal X_{2g,\mathrm{a}}$ on the left hand side of \eqref{EqChangeOfCoor} has the following advantage. The group $P_{2g,\mathrm{a}}(\mathbb R)$ acts transitively on $\mathcal X_{2g,\mathrm{a}}$ by the formula
\[
(v,h) \cdot (v',Z) := (v+hv',hZ)
\]
for $(v,h) \in P_{2g,\mathrm{a}}(\mathbb R)$ and $(v',Z) \in \mathbb R^{2g} \times \mathfrak{H}_g = \mathcal X_{2g,\mathrm{a}}$. The space $\mathfrak{A}_g^{\mathrm{an}}$ is then obtained as the quotient of $\mathcal X_{2g,\mathrm{a}}$ by a congruence subgroup of $P_{2g,\mathrm{a}}(\mathbb Q)$. We refer to \cite[10.5--10.9]{PinkThesis} or \cite[Construction~2.9 and Example~2.12]{Pink05} for these facts.

%We thus obtain the universal covering $\mathcal X_{2g,\mathrm{a}} \rightarrow \mathfrak{A}_g$.

It is clear that both $V_{2g}(\mathbb R)$ and
$\mathrm{Sp}_{2g}(\mathbb R)$ preserve $2\tran{(\rmd a)} \wedge \mathrm{d} b$. Thus this $2$-form is invariant under the action of $P_{2g,\mathrm{a}}(\mathbb R)$ on $\mathcal X_{2g,\mathrm{a}}$.

So by Lemma~\ref{LemmaBettiFormComplexReal}, the previous two paragraphs imply that $\hat{\omega}^{\mathrm{univ}}$ descends to a $(1,1)$-form $\omega^{\mathrm{univ}}$ on $\mathfrak{A}_g$. The semi-positivity of $\omega^{\mathrm{univ}}$ follows from Lemma~\ref{LemmaBettiOnUniversalCovering}.

The property $[N]^*\omega^{\mathrm{univ}} = N^2 \omega^{\mathrm{univ}}$ follows from the ``moreover'' part of Lemma~\ref{LemmaBettiOnUniversalCovering} and the following commutative diagram
\[
\begin{gathered}[b]
\xymatrix{
\IC^g \times \mathfrak{H}_g \ar[r]^-{\widetilde{N}} \ar[d] & \IC^g \times \mathfrak{H}_g \ar[d] \\
\an{\mathfrak{A}}_g \ar[r]^-{[N]} & \an{\mathfrak{A}}_g.
}\\[-\dp\strutbox]
\end{gathered}
\qedhere
\]
\end{proof}


%Thus Proposition~\ref{PropBettiForm} for $\mathfrak{A}_g \rightarrow \mathbb A_g$ follows from Lemma~\ref{LemmaBettiFormUnivAb} and Corollary~\ref{CorBettiFormVanishingDirection}.

This semi-positive $(1,1)$-form $\omega^{\mathrm{univ}}$ will be the Betti form for $\mathfrak{A}_g \rightarrow \mathbb A_g$, as desired in Proposition~\ref{PropBettiForm}. To show this, it suffices to establish property (iii) of Proposition~\ref{PropBettiForm}. Hence it suffices to prove the following proposition. %% Here we use the notations in $\mathsection$\ref{SubsectionBettiMap}: recall the map $b^{\mathrm{univ}} \colon \cA_{\mathfrak{H}_g} \rightarrow \mathbb{T}^{2g}$ in this case and the natural projection $u_B \colon \cA_{\mathfrak{H}_g} = \mathfrak{A}_g \times_{\mathbb A_g} \mathfrak{H}_g \rightarrow \mathfrak{A}_g$.

\begin{proposition}\label{PropBettiFormNonDegUnivAb} 
Assume $\cA \rightarrow S$ is $\mathfrak{A}_g \rightarrow \mathbb A_g$. Let $X$ be an irreducible subvariety of $\mathfrak{A}_g$ of dimension
$d$ and let $\Delta$ be an open subset of $\an{S}$ with $\sman{X}\cap\cA_\Delta\not=\emptyset$. Then
\begin{equation}
\label{prop:bettformvsrankuniv}
 \omega^{\mathrm{univ}}|_{\sman{X}}^{\wedge d} \not\equiv 0
\quad\text{if and only if}\quad
\max_{x \in \sman{X}\cap\cA_\Delta} \mathrm{rank}_{\IR} (\mathrm{d}b_\Delta|_{\sman{X}})_{x} = 2d.
\end{equation}
%where $\widetilde{X}$ is a complex analytic irreducible component of $u_B^{-1}(X^{\mathrm{sm}})$.
\end{proposition}
\begin{proof}
%% Denote by $\widetilde{\omega}^{\mathrm{univ}} := u_B^*(\omega^{\mathrm{univ}})$. 
%% It is clear that
%% \[
%% \omega^{\mathrm{univ}}|_{X^{\mathrm{sm}}}^{\wedge d} \not\equiv 0 \Leftrightarrow \widetilde{\omega}^{\mathrm{univ}}|_{\widetilde{X}}^{\wedge d} \not\equiv 0.
%% \]
%% Thus it suffices to establish the following property. For each $\widetilde{x} \in \widetilde{X}$, we have $(\widetilde{\omega}^{\mathrm{univ}}|_{\widetilde{X}}^{\wedge d})_{\widetilde{x}} = 0$ if and only if
%% \begin{equation}\label{EqRankDegeneracy}
%% \mathrm{rank}_{\mathbb R}(\mathrm{d}b^{\mathrm{univ}}|_{\widetilde{X}})_{\widetilde{x}} < 2d.
%% \end{equation}
We begin by  reformulating
Corollary~\ref{CorBettiFormVanishingDirection}.
If ${C}$ is an irreducible, $1$-dimensional, complex analytic subset
of an open subset of ${\cA_\Delta}$, then
\begin{equation}\label{EqBettiFormVanishingDirection}
{\omega}^{\mathrm{univ}}|_{\sm{C}} = 0 \quad\text{if
and only if}\quad b_\Delta({C})\text{ is a point};
\end{equation}
indeed, this claim is local and it follows using the universal
covering
$\mathbf{u} \colon \IC^g\times\mathfrak{H}_g \rightarrow \mathfrak{A}_g$. 

We assume first that the right side of (\ref{prop:bettformvsrankuniv})
is false, \textit{i.e.}, the maximal rank is strictly less than
$2d=2\dim X$.
So  every $x\in \sman{X}\cap\cA_\Delta$  is a non-isolated point of  
$b_\Delta^{-1}(r) \cap \sman{X}$  where $r =
b_\Delta(x)$. Because 
$b_\Delta^{-1}(r)$ is a complex analytic subset of $\cA_\Delta$ (by Proposition~\ref{PropBettiMap}.(ii) for $\mathfrak{A}_g \rightarrow \mathbb{A}_g$) and $\sman{X}$ is complex
analytic in a neighborhood of $x$ in $\an{\cA}$, there exists  an irreducible
complex analytic curve ${C}$ in  $b_\Delta^{-1}(r) \cap \sman{X}$ passing
through $x$. In
particular, $b_\Delta({C})$ is a point
%and $\widetilde{x} \in \widetilde{C} \subseteq \widetilde{X}$,
and so ${\omega}^{\mathrm{univ}}|_{\sm{C}} \equiv 0$ by \eqref{EqBettiFormVanishingDirection}.

The upshot of the previous paragraph is that the Hermitian form
attached to the semi-positive $(1,1)$-form ${\omega}^{\mathrm{univ}}|_{\sman{X}}$
 vanishes along the tangent space of $\sm{C}$; it is degenerate.
We can complete a tangent vector of $\sm{C}$  to a basis of
the tangent space of $\sman{X}$. Considering holomorphic local
coordinates
we find ${\omega}^{\mathrm{univ}}|_{\sman{X}}^{\wedge d}=0$ at every
point of $\sm{C}$. By continuity it also vanishes at $x \in C$. Since
$x\in \sman{X}\cap\cA_\Delta$ was arbitrary, we conclude
${\omega}^{\mathrm{univ}}|_{\sman{X}}^{\wedge d}\equiv 0$. 

For the converse we assume
${\omega}^{\mathrm{univ}}|_{\sman{X}}^{\wedge d}\equiv 0$. So the
Hermitian form attached to this semi-positive $(1,1)$-form is
degenerate. Thus for each $x \in \sman{X}$, using holomorphic local coordinates we find an irreducible,  $1$-dimensional,  complex analytic subset $C_x$ which passes through $x$ and is contained in $\sman{X}$ such that $\omega^{\mathrm{univ}}|_{\sman{X}}$ vanishes along the tangent space of $\sm{C}_x$. So ${\omega}^{\mathrm{univ}}|_{\sm{C}_x}\equiv 0$, and hence 
%Using holomorphic local coordinates 
% any prescribed auxiliary point of $\sman{X}\cap\cA_\Delta$ lies in an 
%irreducible, $1$-dimensional,  complex analytic subset $C$ of an open
%subset of  $\sman{X}$ with
%${\omega}^{\mathrm{univ}}|_{\sm{C}}\equiv 0$. So
 $b_\Delta(C_x)$ is a point
by \eqref{EqBettiFormVanishingDirection}.
Letting the point $x$ run over $\sman{X}$, we conclude that the rank on the right
side of (\ref{prop:bettformvsrankuniv}) is strictly less than $2d$. 
\end{proof}


\subsection{General case}
We now prove Propositions~\ref{PropBettiMap}
and \ref{PropBettiForm} for
 $\pi \colon \cA \rightarrow S$ as near the beginning of this section.
 In particular, we assume \texttt{(Hyp)}. With the construction in $\mathsection$\ref{SubsectionBettiMap}, the rest of the proof of Proposition~\ref{PropBettiMap} follows the construction in \cite[$\mathsection$4]{GaoBettiRank}. 

%% By assumption, $\cA$  carries a principal polarization
%% and  level-$\ell$-structure.

As
$\IA_g$ is a fine moduli space there exists a Cartesian diagram
\begin{equation*}%\label{EqEmbedAbSchIntoUnivAbVar}
\xymatrix{
\cA \ar[r]^-{\iota} \ar[d]_{\pi} \pullbackcorner & \mathfrak{A}_g \ar[d] \\
S \ar[r]^-{\iota_S} & \mathbb{A}_g.
}
\end{equation*}

Now let $s_0 \in S(\IC)$. Applying Proposition~\ref{PropBettiMap} to the universal abelian variety $\mathfrak{A}_g \rightarrow \mathbb A_g$ and $\iota_S(s_0) \in \mathbb A_g(\IC)$, we obtain an open neighborhood $\Delta_0$ of $\iota_S(s_0)$ in $\mathbb A_g^{\mathrm{an}}$ and a map
\[
b_{\Delta_0} \colon \mathfrak{A}_g|_{\Delta_0} \rightarrow \mathbb T^{2g}
\]
satisfying the properties listed in Proposition~\ref{PropBettiMap}.

Now let $\Delta = \iota_S^{-1}(\Delta_0)$. Then $\Delta$ is an open neighborhood of $s$ in $S^{\mathrm{an}}$. Denote by $\cA_{\Delta} = \pi^{-1}(\Delta)$ and define
\[
b_{\Delta} = b_{\Delta_0} \circ \iota \colon \cA_{\Delta} \rightarrow \mathbb T^{2g}.
\]
Then $b_{\Delta}$ satisfies the properties listed in Proposition~\ref{PropBettiMap} for $\cA \rightarrow S$. Hence $b_{\Delta}$ is our desired Betti map.


Next let us turn to the Betti form. Let $\omega^{\mathrm{univ}}$ be the semi-positive $(1,1)$-form on $\mathfrak{A}_g$ as in Lemma~\ref{LemmaBettiFormUnivAb}. Define $\omega:= \iota^*\omega^{\mathrm{univ}}$. We will show that $\omega$ satisfies the properties listed in Proposition~\ref{PropBettiForm}.

The $(1,1)$-form $\omega$ is semi-positive as it is the pull-back of
the semi-positive form $\omega^{\mathrm{univ}}$. Moreover, it satisfies
$[N]^*\omega = N^2 \omega$ since $\omega^{\mathrm{univ}}$ has this
property. Hence we have established properties (i) and (ii) of
Proposition~\ref{PropBettiForm}.

Let us verify (iii) of Proposition~\ref{PropBettiForm}. Suppose $X$ is
an irreducible subvariety of $\cA$ of dimension $d$.
Let $\Delta$ be an open subset of $\an{S}$ with
$\sman{X}\cap\cA_\Delta\not=\emptyset$; we may shrink $\Delta$ subject
to this condition. 
Let $Z
= \overline{\iota(X)}$ and observe $\dim Z \le d$. 

Since $\omega = \iota^*\omega^{\mathrm{univ}}$, we have
\begin{equation}\label{EqOmegaAndOmegaUniv}
\omega|_{\sman{X}}^{\wedge d} \not\equiv 0 \quad\text{if and only if}\quad \omega^{\mathrm{univ}}|_{\sman{Z}}^{\wedge d} \not\equiv 0.
\end{equation}
Next by definition of $b_{\Delta}$, we have the following property:
 For suitable non-empty open subsets $\Delta$ of $S^{\mathrm{an}}$ and
 $\Delta_0$ of $\mathbb{A}_g^{\mathrm{an}}$
 %$\overline{\iota_S(S)}^{\mathrm{an}}$
 such that $\iota_S(\Delta) \subseteq \Delta_0$, we have
\begin{equation}\label{EqBettiRankTrivialBound}
\max_{x \in \sman{X} \cap \cA_{\Delta}} \mathrm{rank}_{\IR}
(\mathrm{d}b_{\Delta}|_{\sman{X}})_x
\le \max_{x \in \sman{Z} \cap \mathfrak{A}_{g,\Delta_0}} \mathrm{rank}_{\IR} (\mathrm{d}b_{\Delta_0}|_{Z^{\mathrm{sm,an}}})_x \le 2\dim Z \le 2d.
\end{equation}

Suppose first that $\omega|_{\sman{X}}^{\wedge d}\not\equiv 0$, then
\eqref{EqOmegaAndOmegaUniv} implies
$\omega^{\mathrm{univ}}|_{\sman{Z}}^{\wedge d}\not\equiv 0$ and in
particular $d=\dim Z$. 
We can apply 
Proposition~\ref{PropBettiForm}(iii)  to $Z$ and obtain
$ \max_{x \in \sman{Z} \cap \mathfrak{A}_{g,\Delta_0}} \mathrm{rank}_{\IR}
(\mathrm{d}b_{\Delta_0}|_{Z^{\mathrm{sm,an}}})_x =2d$. 
Now $\iota|_X\colon X\rightarrow Z$ is
generically finite as $\dim X = \dim Z$,
so the first inequality in \eqref{EqBettiRankTrivialBound} is an
equality. We conclude
\begin{equation}
\label{eq:maxrankXgeneralA}
\max_{x \in \sman{X} \cap \cA_\Delta} \mathrm{rank}_{\IR}
(\mathrm{d}b_{\Delta}|_{X^{\mathrm{sm,an}}})_x =2d.
\end{equation}


Conversely, assume \eqref{eq:maxrankXgeneralA} holds true.
Then we have equalities throughout in \eqref{EqBettiRankTrivialBound}.
By Proposition~\ref{PropBettiForm}(iii) applied to $Z$ and
by \eqref{EqOmegaAndOmegaUniv} we get $\omega|_{\sman{X}}^{\wedge
d}\not\equiv 0$. 




%%% Local Variables:
%%% TeX-master: "main"
%%% End:
